\section{Discussion}
\label{sec:discussion}

The central-path Hessian is constrained by an object that does not
refer to the barrier: the difference body of the current objective
sublevel. This explains both the matching condition-number law and
the comparison of complete spectra at equal accuracy. Within a fixed
representation and metric, a finite-parameter barrier change cannot
remove a conditioning exponent imposed by sublevel geometry. Uniform
bounds on barrier parameters are essential if the barrier itself varies
with the requested gap.

The resulting geometric distinction is finer than several customary
classifications. A unique optimum does not by itself imply bounded
conditioning outside the polyhedral setting. Conversely, primal
vertex degeneracy does not prevent a positive-definite scaled
log-Hessian limit in an LP. In SDP, an attained fractional error bound
can determine the conditioning exponent, but singularity degree alone
cannot: the paired degree-two systems have different diameter rates.
Full-spectrum width bounds explain why an extreme-eigenvalue ratio
also need not describe the interior spectral scales.

For linear solves, the direction of the forcing and the accuracy
contract remain decisive. On an LP central tail, the weak eigenspace
approaches the optimal-face tangent and the exact Newton forcing has
a small projection onto it. The sharp barrier example shows why this
fact is insufficient for accurate direction recovery after discarding
weak modes. Clustered CG can nevertheless exploit separated spectral
intervals in exact arithmetic. Its conclusion is an energy-error
bound, and the numerical examples expose the need to distinguish it
from a recursively computed residual. Quantum applications require
additional access and output assumptions before either fact can be
converted into a runtime statement.

These are statements about a specified primal operator. They coexist
with useful changes of representation, gap-dependent preconditioners,
and primal--dual formulations whose spectra can behave differently.
The precise comparison established here supplies a common geometric
reference for assessing such changes: it identifies what is forced
by objective sublevels before accounting for the cost and effect of
the linear algebra used to solve the Newton equations.
